%=============================================================================
% WAVE 3 ITERATION 6: Patent 8 (Communications) + Patent 11 (EM Coupling)
%
% DFD Research Programme -- March 2026
% Agent 8/11, Iteration 6 of 20
%
% W3i5 RESULTS INHERITED:
%   P8: M_eff zero impact; DSN wins 370x; P15 transmitter killed
%   P11: 11 configurations beat SQMS; bulk SQUID hopeless; MEMS path viable
%   Cosmo: Two-component psi = psi_grav + delta_psi_EM
%   P2: Passive Superiority Theorem; Volume Cancellation Theorem
%
% W3i6 MANDATE:
%   P8: Under two-component decomposition, psi-communications use
%       delta_psi_EM (EM-sourced). Does inherent secrecy change?
%       Does corridor propagation change? Honest: what is the ONE
%       killer application where psi-comms beats ALL alternatives?
%   P11: With 11 viable configurations identified, RANK them by:
%       (a) cost, (b) timeline, (c) |lambda-1| reach, (d) systematic risk.
%       Produce the definitive detection roadmap.
%       Which experiment should be built FIRST?
%=============================================================================

\documentclass[11pt]{article}
\usepackage[margin=2cm]{geometry}
\usepackage{amsmath,amssymb,amsthm,mathtools}
\usepackage{physics}
\usepackage{booktabs}
\usepackage{array}
\usepackage{longtable}
\usepackage{hyperref}
\usepackage{xcolor}
\usepackage{siunitx}
\usepackage{bm}
\usepackage{enumitem}
\usepackage{tcolorbox}
\usepackage{float}

\newtheorem{theorem}{Theorem}[section]
\newtheorem{proposition}[theorem]{Proposition}
\newtheorem{corollary}[theorem]{Corollary}
\newtheorem{lemma}[theorem]{Lemma}
\theoremstyle{definition}
\newtheorem{definition}[theorem]{Definition}
\theoremstyle{remark}
\newtheorem{remark}[theorem]{Remark}
\newtheorem{finding}[theorem]{Finding}
\newtheorem{correction}[theorem]{Correction}

\newcommand{\ee}[1]{\times 10^{#1}}
\newcommand{\psifield}{\psi}
\newcommand{\psigrav}{\psi_{\rm grav}}
\newcommand{\dpsiEM}{\delta\psi_{\rm EM}}
\newcommand{\lam}{\lambda}
\newcommand{\lamdiff}{|\lambda-1|}
\newcommand{\Meff}{M_\psi^{\rm eff}}
\newcommand{\Mgrav}{M_\psi^{\rm grav}}
\newcommand{\muG}{\mu_0^{\rm gal}}
\newcommand{\Vcav}{V_{\rm cav}}
\newcommand{\Opsi}{\Omega_\psi}
\newcommand{\SNR}{\mathrm{SNR}}
\newcommand{\BER}{\mathrm{BER}}
\newcommand{\Cs}{C_{\rm s}}
\newcommand{\DFD}{\textsc{dfd}}
\newcommand{\GN}{G_N}
\newcommand{\dpsi}{\delta\psi}
\DeclareMathOperator{\Tr}{Tr}

\newcommand{\confirmed}[1]{\textcolor{blue}{\textbf{CONFIRMED:} #1}}
\newcommand{\corrected}[1]{\textcolor{red}{\textbf{CORRECTED:} #1}}
\newcommand{\caution}[1]{\textcolor{orange}{\textbf{CAUTION:} #1}}
\newcommand{\honest}[1]{\textcolor{violet}{\textbf{HONEST ASSESSMENT:} #1}}
\newcommand{\killed}[1]{\textcolor{red!70!black}{\textbf{KILLED:} #1}}
\newcommand{\verdict}[1]{\textcolor{green!50!black}{\textbf{VERDICT:} #1}}

\tcbuselibrary{skins,breakable}

\title{\textbf{Wave 3 Iteration 6: Cross-Reference Update}\\[6pt]
Patent 8 --- Communications\\
Patent 11 --- EM Coupling, SRF Cavity, Psi-Detection\\[4pt]
{\normalsize Two-Component $\psi$ Impact on Secrecy and Propagation;}\\
{\normalsize The One Killer Application for Psi-Comms;}\\
{\normalsize Definitive Detection Roadmap: 11 Configurations Ranked}}
\author{DFD Research Programme --- P8/P11 Combined Agent, Iteration 6}
\date{20 March 2026}

\begin{document}
\maketitle
\tableofcontents
\newpage

%=============================================================================
\section*{Executive Summary: W3i6 P8/P11}
%=============================================================================

\begin{tcolorbox}[colback=blue!5!white, colframe=blue!60!black,
  title=CRITICAL W3i6 RESULTS --- READ FIRST]
\begin{enumerate}[nosep]

\item \textbf{P8: Two-component $\psi$ CLARIFIES communications physics.}
  Under the W3i5 decomposition $\psi = \psigrav + \dpsiEM$, psi-communications
  operates on $\dpsiEM$ (EM-sourced perturbation), not on $\psigrav$
  (gravitational background). This is a structural clarification, not a
  correction: the W3i4/W3i5 link budget used $\delta\psi \propto GM_s/(c^2 d)$,
  which is correct for modulating a macroscopic mass (the modulation
  couples through the equivalence principle regardless of source decomposition).
  \textbf{Secrecy is unchanged.} \textbf{Corridor propagation is unchanged.}

\item \textbf{P8: The ONE killer application is submarine/underground
  communication.} After honest elimination of every scenario where EM
  alternatives exist, the single use case where psi-comms is irreplaceable
  is \textbf{bidirectional real-time communication with deeply submerged
  submarines or underground facilities}. The psi-field propagates through
  matter with absorption coefficient $\alpha < 10^{-20}$~m$^{-1}$ (W3i4),
  compared with EM skin-depth attenuation of $>1$~dB/m at useful
  frequencies. No EM, acoustic, or neutrino alternative provides
  real-time bidirectional communication at depth.

\item \textbf{P8: HONEST CAVEAT.} Even the killer application requires a
  modulated mass source of $M_s \gtrsim 10^{10}$~kg for useful data rates
  at submarine operating depths, which remains technologically infeasible.
  The psi-channel is theoretically unique but practically unrealisable
  with any foreseeable technology.

\item \textbf{P11: DEFINITIVE DETECTION ROADMAP PRODUCED.} All 11
  configurations ranked by cost, timeline, $\lamdiff$ reach, and systematic
  risk. The \textbf{first experiment to build} is an SQMS cavity array
  (4 cavities, incoherent averaging) at Fermilab, reaching
  $\lamdiff \sim 7.8\ee{-10}$ in 2027--2028 at cost $\sim$\$2--4M.
  This is the minimum viable experiment that improves on the current
  SQMS bound with established technology.

\item \textbf{P11: Three-phase roadmap.}
  Phase~I (2027--2028): SQMS 4-cavity array, $\lamdiff < 7.8\ee{-10}$.
  Phase~II (2028--2030): P2+P11 hybrid with $\mu$g MEMS, $\lamdiff < 4.7\ee{-11}$.
  Phase~III (2030--2035): 256-cavity coherent array or P2+P11 with ng MEMS,
  $\lamdiff < 10^{-12}$.

\end{enumerate}
\end{tcolorbox}

%=============================================================================
%                    PART I: PATENT 8 (COMMUNICATIONS)
%=============================================================================

\part{Patent 8 --- Communications Under Two-Component $\psi$}

%=============================================================================
\section{The Two-Component Decomposition: Impact on P8 Physics}
\label{sec:two_component_p8}
%=============================================================================

\subsection{Recall: The W3i5 Two-Component $\psi$-Field}

W3i5 (Cosmo/MatSci update) established the structural resolution of tensions
T1, T2, and T5:
\begin{equation}
\boxed{
\psi = \psigrav + \dpsiEM,
}
\label{eq:two_component}
\end{equation}
where:
\begin{itemize}[nosep]
\item $\psigrav$ is the gravitational potential, sourced by all stress-energy
  through the metric identification $g_{00} = -e^{2\psi}$. Numerically,
  $\psigrav^{\rm cosmo} \approx 0.047$ (Mach integral).
\item $\dpsiEM$ is the EM-sourced perturbation governed by the DFD field
  equation $\Box\psi + (1+\kappa)[(\nabla\psi)^2 - \dot\psi^2/c^2]
  = 4\pi \GN T_{\rm EM}/c^4$. Numerically,
  $\dpsiEM^{\rm cosmo} \approx 9.6\ee{-6}$.
\end{itemize}

\subsection{Which Component Does P8 Use?}

\begin{finding}[P8 Communications Uses Both Components]
\label{find:p8_both_components}
The P8 psi-channel operates by \emph{modulating the gravitational field
of a macroscopic mass}. When a mass $M_s$ is modulated at frequency $f_c$
with depth $\xi_0$, the resulting psi-field perturbation at distance $d$ is:
\begin{equation}
\delta\psi_{\rm signal}(t) = \frac{G M_s \xi_0}{c^2 d}\cos(2\pi f_c t).
\label{eq:psi_signal}
\end{equation}
This perturbation arises from $\psigrav$ --- it is a modulation of the
Newtonian gravitational potential, sourced by the mass redistribution.

However, the DFD field equation states that $\psi$-perturbations propagate
with the DFD dispersion relation (speed $c_s = c/n_\psi$, where
$n_\psi = 1$ in the Newtonian regime). The propagation characteristics
are determined by the \emph{kinetic sector} of the field equation, which
contains the $\mu(|\nabla\psi|/a_0)$ interpolation function. This kinetic
sector governs both $\psigrav$ and $\dpsiEM$ perturbations.

\textbf{Key insight}: In the two-component picture, the signal
$\delta\psi_{\rm signal}$ from mass modulation is a perturbation of
$\psigrav$. But the same field equation governs its propagation.
The coupling constant $(\lam-1)$ enters only when EM energy is
\emph{directly} converted to psi-field energy (Process~A). For P8
mass-modulation communications, the signal is generated gravitationally,
not electromagnetically.

\textbf{Consequence}: The two-component decomposition does not change any
P8 prediction. The link budget (Eq.~\ref{eq:psi_signal}) uses $\psigrav$
physics. The receiver detects gravitational psi-perturbations. The
$(\lam-1)$ coupling is relevant only if one attempts P15-type EM-to-psi
transmission (which was killed in W3i5).
\end{finding}

%=============================================================================
\section{Does Inherent Secrecy Change Under Two-Component $\psi$?}
\label{sec:secrecy_two_component}
%=============================================================================

\subsection{Secrecy Mechanism Audit}

The W3i4/W3i5 secrecy capacity is:
\begin{equation}
\Cs = B\log_2\!\left(\frac{1 + \SNR_B}{1 + \SNR_E}\right) > 0
\quad\text{iff}\quad \SNR_B > \SNR_E.
\label{eq:secrecy_capacity}
\end{equation}
The inequality $\SNR_B > \SNR_E$ follows from:
\begin{enumerate}
\item \textbf{Corridor geometry}: Bob inside the 6-observable corridor
  receives MRC-combined signal; Eve outside receives only the scalar
  monopole (far-field).
\item \textbf{Hadamard bound}: Eve's mutual information is exponentially
  suppressed with distance from the corridor.
\end{enumerate}

\begin{theorem}[Secrecy Under Two-Component $\psi$: Unchanged]
\label{thm:secrecy_unchanged}
The inherent secrecy of the psi-channel is determined by the
\emph{tensor structure} of the psi-field perturbation. A mass modulation
generates perturbations in all six observables
($\psi, \partial_i\psi, \partial_i\partial_j\psi$), with the
higher-order observables decaying faster than the monopole ($1/d^2$
vs $1/d$). The corridor advantage arises because Bob can measure all
six components while Eve (at large distance) measures only the monopole.

This tensor structure is a property of the \emph{gravitational}
psi-field ($\psigrav$) and its derivatives. Under the two-component
decomposition:
\begin{enumerate}[label=(\roman*)]
\item The signal is a $\psigrav$ perturbation --- tensor structure unchanged.
\item The noise floor is set by the receiver technology (atom
  interferometer), not by $\dpsiEM$ --- noise unchanged.
\item The corridor geometry is a consequence of $1/r^n$ decay of
  multipole moments --- geometry unchanged.
\item Eve's Hadamard bound depends only on the spatial correlation
  length of the field --- unchanged.
\end{enumerate}

\confirmed{Inherent secrecy is \textbf{completely unaffected} by the
two-component decomposition. The secrecy mechanism operates in the
$\psigrav$ sector, which is the standard gravitational sector.}
\end{theorem}

\subsection{A Subtlety: Does $\dpsiEM$ Create a Noise Floor?}

One might worry that the cosmological $\dpsiEM$ background
($\dpsiEM^{\rm cosmo} \approx 9.6\ee{-6}$) introduces additional noise
in the psi-channel. However:
\begin{itemize}[nosep]
\item The \emph{fluctuations} of $\dpsiEM$ at laboratory frequencies
  ($f \sim 1$~kHz) are negligible: the cosmological EM background (CMB,
  starlight) is static on communication timescales.
\item Local EM sources (laboratory equipment, power lines) produce
  $\dpsiEM$ perturbations at the $(\lam-1)$ level, which is
  $< 1.56\ee{-9}$ of their EM amplitude. This is far below the
  gravitational psi-noise floor.
\item The dominant noise source remains the quantum projection noise of
  the atom interferometer receiver, not any $\dpsiEM$ contamination.
\end{itemize}

\begin{finding}[$\dpsiEM$ Noise Negligible for P8]
\label{find:dpsiEM_noise}
The EM-sourced component $\dpsiEM$ does not contribute a meaningful noise
floor to psi-communications. The $(\lam-1)$ suppression ensures that any
EM noise leaking through Process~A is at least $10^{9}\times$ below the
gravitational signal level. The P8 noise budget (W3i4 Table~3) is
\textbf{unchanged}.
\end{finding}

%=============================================================================
\section{Does Corridor Propagation Change?}
\label{sec:corridor_propagation}
%=============================================================================

\subsection{Propagation Speed}

The psi-field propagation speed in the two-component picture:
\begin{itemize}[nosep]
\item $\psigrav$ perturbations propagate at $c_s = c/n_\psi$, where
  $n_\psi = 1$ in the Newtonian regime ($|\nabla\psi|/a_0 \gg 1$) and
  $n_\psi = \sqrt{2}$ in the MOND regime ($|\nabla\psi|/a_0 \ll 1$).
\item $\dpsiEM$ perturbations are governed by the same kinetic sector
  of the field equation and propagate at the same speed $c_s$.
\end{itemize}

For P8 communications (mass modulation generating $\psigrav$ perturbations):
\begin{equation}
v_{\rm prop} = c/n_\psi \approx c \quad\text{(surface/near-surface operations)}.
\end{equation}
This is unchanged by the two-component decomposition.

\subsection{Corridor Formation}

The corridor is defined by the region where the receiver can detect
all six psi-observables with $\SNR > 1$. The corridor geometry is set by
the source mass distribution and the $1/r^n$ multipole decay. Under the
two-component picture:

\begin{finding}[Corridor Propagation Unchanged]
\label{find:corridor_unchanged}
\confirmed{The corridor geometry, propagation speed, and signal
structure are all determined by the $\psigrav$ sector and are
\textbf{completely unaffected} by the two-component decomposition.}
The corridor width $w \sim \sqrt{\lambda_\psi d}$ (where
$\lambda_\psi = c_s/f_c$ is the psi-wavelength) and the corridor
length (limited by SNR degradation with distance) are unchanged from
W3i4/W3i5 values.
\end{finding}

%=============================================================================
\section{Honest Assessment: What Psi-Comms Beats and What It Does Not}
\label{sec:honest_assessment}
%=============================================================================

\subsection{Systematic Elimination of Candidate Applications}

We now perform a rigorous elimination of every proposed psi-comms
application to identify where (if anywhere) it is genuinely
irreplaceable.

\begin{longtable}{@{}p{3.8cm}p{3.5cm}p{3.5cm}p{3cm}@{}}
\caption{Systematic comparison: psi-comms vs best EM/alternative for
each proposed application. DFD numbers assume validity of DFD physics
and availability of required source masses.}
\label{tab:application_elimination}\\
\toprule
Application & Psi-Comms & Best Alternative & Verdict \\
\midrule
\endfirsthead
\toprule
Application & Psi-Comms & Best Alternative & Verdict \\
\midrule
\endhead
\bottomrule
\endlastfoot
%
Earth--Mars opposition & 14 kbits/s (requires $10^{24}$~kg TX) &
  DSN: 5.2 Mbits/s (35~kg HGA) & \killed{DSN $370\times$ faster, $10^{22}\times$ lighter} \\
\midrule
Earth--Mars conjunction & 8.7 kbits/s (no blackout) &
  DSN: $\sim 0$ (2-week blackout) &
  \caution{DFD wins on availability; DSN rate is higher outside blackout} \\
\midrule
Outer solar system ($>100$~AU) & $\sim 130$ bits/s at 150~AU
  (requires $10^{24}$~kg TX) &
  Laser comms: $\sim 1$~kbits/s projected (DSOC) &
  \killed{Laser comms projected to win; TX mass absurd} \\
\midrule
Interstellar ($>1$~pc) & $C \propto \log(1/d^2)$ advantage &
  EM: $C \propto 1/d^2$ &
  \caution{DFD theoretically superior but requires stellar-mass TX} \\
\midrule
Submarine at depth & Zero attenuation through seawater &
  VLF/ELF: 1-way only, $\sim 1$~bit/min; acoustic: $<10$~km range &
  \verdict{DFD irreplaceable for bidirectional deep-submarine comms} \\
\midrule
Underground (mine/bunker) & Zero attenuation through rock &
  Through-the-earth (TTE): $<100$~m practical depth &
  \verdict{DFD irreplaceable for deep underground comms} \\
\midrule
Secure military comms & Inherent $C_s > 0$ (no key) &
  QKD + AES-256: proven, deployed &
  \killed{QKD + symmetric encryption is adequate; DFD TX mass absurd} \\
\midrule
Disaster/earthquake comms & Propagation through rubble &
  Satellite relay: requires clear sky &
  \caution{DFD niche advantage for buried survivors; TX mass problem persists} \\
\bottomrule
\end{longtable}

\subsection{Elimination Reasoning}

\begin{enumerate}
\item \textbf{Earth--Mars}: Eliminated. DSN wins on rate by $370\times$ and
  on TX mass by $10^{22}\times$. The conjunction blackout advantage is real
  but marginal: conjunction lasts $\sim 2$ weeks in a 26-month synodic cycle
  ($<2\%$ of time), and mission planning accommodates this routinely.
  Laser communications (DSOC) will further extend EM capability.

\item \textbf{Outer solar system}: Eliminated. The DSOC demonstration
  (2023--2024) achieved 267~Mbits/s from 0.1~AU. Projected deep-space laser
  relay networks will reach $>1$~kbits/s at 100~AU with $<100$~kg terminals.
  DFD requires Venus-mass transmitters.

\item \textbf{Interstellar}: Theoretically favourable ($\log$ vs $1/d^2$
  scaling) but requires modulating a solar-mass object. Eliminated on
  feasibility.

\item \textbf{Secure military}: Eliminated. Quantum key distribution (QKD)
  combined with one-time-pad or AES-256 provides information-theoretic
  security without the $10^{24}$~kg transmitter mass problem. DFD's
  inherent secrecy is elegant but impractical.

\item \textbf{Submarine}: \textbf{NOT eliminated.} Current submarine
  communications use VLF (3--30~kHz) with one-way, very-low-rate ($<100$
  bits/s) shore-to-sub only, requiring trailing wire antennas near the
  surface. ELF (3--300~Hz) penetrates deeper but at $\sim 1$~bit/min.
  There is \emph{no existing or projected technology} for real-time
  bidirectional communication with a submarine at operating depth
  ($>100$~m). The psi-channel provides:
  \begin{itemize}[nosep]
  \item Zero attenuation through seawater ($\alpha < 10^{-20}$~m$^{-1}$).
  \item Bidirectional communication (sub can modulate its own mass
    distribution for uplink).
  \item Depth-independent: works at $100$~m or $1000$~m.
  \end{itemize}

\item \textbf{Underground}: \textbf{NOT eliminated.} Through-the-earth
  (TTE) communications for mining emergencies achieves $<100$~m practical
  depth with extremely low rates. No technology penetrates kilometres
  of rock at useful data rates. DFD psi-comms would be unique.

\item \textbf{Disaster/earthquake}: Partially surviving. The ability
  to communicate through arbitrary amounts of rubble/debris is unique,
  but the TX mass problem applies. Eliminated as primary application.
\end{enumerate}

%=============================================================================
\section{The ONE Killer Application}
\label{sec:killer_app}
%=============================================================================

\begin{tcolorbox}[colback=green!5!white, colframe=green!50!black,
  title=\textbf{THE ONE KILLER APPLICATION: Submarine/Underground Communication}]

\textbf{Application}: Bidirectional, real-time communication with deeply
submerged submarines or underground facilities through arbitrary thicknesses
of seawater, rock, or soil.

\textbf{Why psi-comms is irreplaceable}: The psi-field absorption
coefficient in matter is $\alpha_\psi < 10^{-20}$~m$^{-1}$ (W3i4
Sec.~VII). For comparison:
\begin{center}
\begin{tabular}{lll}
\toprule
Medium & EM skin depth (1~kHz) & Psi absorption length \\
\midrule
Seawater & $\sim 8$~m & $>10^{20}$~m \\
Granite & $\sim 50$~m & $>10^{20}$~m \\
Basalt & $\sim 30$~m & $>10^{20}$~m \\
Iron ore & $\sim 5$~m & $>10^{20}$~m \\
\bottomrule
\end{tabular}
\end{center}

No EM frequency penetrates more than a few skin depths. Acoustic
communication in seawater is limited to $\sim 10$~km range at low
rates. Neutrino communication (proposed) requires $\sim$~PW accelerator
power for $\sim 0.1$~bits/s.

\textbf{The psi-channel is the ONLY known physics that permits
bidirectional communication through arbitrary matter thickness at
useful data rates.}

\textbf{Honest caveat --- the transmitter mass problem}: Even for the
submarine application, the data rate scales as $\log_2(1 + G^2 M_s^2
\xi_0^2/(c^4 d^2 N_0 B))$. For a submarine at $d = 500$~m depth, the
required source mass for 1~kbit/s is:
\begin{equation}
M_s \approx \frac{c^2 d}{G\xi_0}\sqrt{\frac{(2^{C/B}-1)N_0 B}{\eta}}
\sim 10^{10}\text{~kg} \quad\text{(at $d = 500$~m, $C = 1$~kbit/s)}.
\end{equation}
This is $\sim 10^7$ tonnes --- roughly the mass of a large asteroid
fragment. For a shore-based installation with a sufficiently massive
foundation (e.g., modulating the mass distribution of a mountain or
reservoir), this is conceivable in principle but not with current
engineering.

\textbf{Reduced-rate scenario}: At $M_s = 10^6$~kg (1000 tonnes,
achievable with engineered mass redistribution):
\begin{equation}
\SNR = \frac{G^2 (10^6)^2 (10^{-3})^2}{c^4 (500)^2 \times 2.5\ee{-38}
\times 10^3} \approx \frac{4.4\ee{-22} \times 10^{-6}}{8.1\ee{33}
\times 2.5\ee{-35} \times 10^3} = \frac{4.4\ee{-28}}{2.0\ee{2}}
\approx 2.2\ee{-30}.
\end{equation}
This gives effectively zero rate. The gap between ``technologically
achievable'' transmitter masses ($\lesssim 10^6$~kg) and the required
mass ($\gtrsim 10^{10}$~kg) spans four orders of magnitude.

\killed{Even the killer application is currently technologically
infeasible. The transmitter mass gap of $\sim 10^4\times$ cannot be
bridged by any foreseeable engineering.}

\textbf{However}: this is the unique application where \emph{no
alternative physics} can compete. If a breakthrough in mass-modulation
technology (or a reduction in the required $M_s$ through coherent
enhancement or P2 passive amplification) were achieved, submarine
psi-comms would be transformative.
\end{tcolorbox}

%=============================================================================
\section{P8 Summary: Status After W3i6}
\label{sec:p8_summary}
%=============================================================================

\begin{table}[h]
\centering
\caption{P8 claim status after W3i6 two-component audit.}
\begin{tabular}{@{}p{4.5cm}p{2.5cm}p{6cm}@{}}
\toprule
Claim & Status & Note \\
\midrule
Two-component changes link budget & \confirmed{No change} &
  Signal is $\psigrav$; $\dpsiEM$ irrelevant \\
Secrecy under two-component & \confirmed{Unchanged} &
  Tensor structure in $\psigrav$ sector \\
Corridor propagation & \confirmed{Unchanged} &
  Kinetic sector governs both components \\
$\dpsiEM$ noise floor & \confirmed{Negligible} &
  $(\lam-1)$ suppression $>10^9\times$ \\
Deep-space comms (Mars+) & \killed{Infeasible} &
  DSN $370\times$ faster, TX mass absurd \\
Submarine/underground & \verdict{Unique but infeasible} &
  Only physics that works; TX mass gap $10^4\times$ \\
ONE killer application & Submarine/underground & No alternative physics; massive
  engineering gap \\
\bottomrule
\end{tabular}
\label{tab:p8_summary_w3i6}
\end{table}

%=============================================================================
%                    PART II: PATENT 11 (EM COUPLING)
%=============================================================================

\newpage
\part{Patent 11 --- Definitive Detection Roadmap}

%=============================================================================
\section{Inventory: The 11 Configurations That Beat SQMS}
\label{sec:inventory}
%=============================================================================

W3i5 identified 11 configurations with $\lamdiff < 1.56\ee{-9}$ (the
SQMS bound). We now assign to each: (a)~estimated cost, (b)~timeline
to first result, (c)~$\lamdiff$ reach, and (d)~dominant systematic risk.

\subsection{Cost and Timeline Methodology}

Cost estimates are based on:
\begin{itemize}[nosep]
\item SRF cavity costs from Fermilab SQMS programme ($\sim$\$200k per
  TESLA 9-cell cavity, dressed).
\item MEMS fabrication costs from published DARPA/NIST programmes.
\item SQUID electronics at $\sim$\$50k per readout channel.
\item Cryogenic infrastructure: existing (Fermilab, DESY, ESS) vs
  purpose-built.
\item Personnel at \$250k FTE/yr (fully loaded).
\end{itemize}

Timeline estimates assume:
\begin{itemize}[nosep]
\item Phase~I: uses existing infrastructure and technology.
\item Phase~II: requires one new technology development (MEMS integration).
\item Phase~III: requires either scaling (256 cavities) or frontier
  technology (sub-ng MEMS).
\end{itemize}

%=============================================================================
\section{The Definitive Ranking}
\label{sec:ranking}
%=============================================================================

\begin{longtable}{@{}p{0.5cm}p{3.2cm}p{1.5cm}p{1.5cm}p{1.5cm}p{4.5cm}@{}}
\caption{All 11 sub-SQMS configurations ranked by composite score
(cost $\times$ timeline / reach). Lower composite = higher priority.
Cost in \$M, timeline in years from approval, reach as $\lamdiff$.}
\label{tab:definitive_ranking}\\
\toprule
\# & Configuration & Cost (\$M) & Timeline (yr) & $\lamdiff$ Reach &
  Dominant Systematic Risk \\
\midrule
\endfirsthead
\toprule
\# & Configuration & Cost (\$M) & Timeline (yr) & $\lamdiff$ Reach &
  Dominant Systematic Risk \\
\midrule
\endhead
\bottomrule
\endlastfoot
%
1 & SQMS 4-cavity (incoh.) & 2--4 & 1--2 & $7.8\ee{-10}$ &
  Correlated systematics (vibration, magnetic) across cavities \\
\midrule
2 & 256-cavity (incoh.) & 60--80 & 3--4 & $9.75\ee{-11}$ &
  Correlated systematics at $\sqrt{N}$ limit; infrastructure scale \\
\midrule
3 & 256-cavity (coherent) & 80--120 & 4--5 & $6.1\ee{-12}$ &
  Phase coherence across 256 cavities at GHz; no demonstrated
  psi-phased-array \\
\midrule
4 & P2+P11 hybrid, $\mu$g MEMS & 3--6 & 2--3 & $4.7\ee{-11}$ &
  MEMS--SRF integration at cryogenic; mechanical $Q$ in SRF
  environment \\
\midrule
5 & P2+P11 hybrid, ng MEMS & 5--10 & 3--4 & $1.5\ee{-12}$ &
  ng-scale MEMS with $Q_{\rm mech} = 10^6$ undemonstrated at 2~K \\
\midrule
6 & P2+P11 hybrid, pg MEMS & 10--20 & 5--7 & $4.7\ee{-14}$ &
  pg resonators are research frontier; integration with SRF
  speculative \\
\midrule
7 & P2+P11 hybrid, fg MEMS & 20--50 & 7--10+ & $1.5\ee{-15}$ &
  fg mechanical resonators do not yet exist at required $Q$ \\
\midrule
8 & P2+P11 hybrid, ag MEMS & 50--100 & 10+ & $4.7\ee{-17}$ &
  Levitated nanoparticles at mK; extreme research frontier \\
\midrule
9 & ESS Lund (realistic) & 5--15 & 3--5 & $10^{-12}$--$10^{-13}$ &
  Requires ESS beam time allocation; neutron cross-correlation
  systematics \\
\midrule
10 & ESS Lund (optimistic) & 5--15 & 3--5 & $1.14\ee{-14}$ &
  Assumes systematic control at $10^{-14}$ level; may be
  over-optimistic \\
\midrule
11 & LIGO 6th channel & 0.5--2 & 2--3 & $6\ee{-19}$ &
  Extraction of psi-signal from LIGO data stream; enormous
  backgrounds; no demonstrated methodology \\
\end{longtable}

%=============================================================================
\section{Detailed Assessment of Each Configuration}
\label{sec:detailed_assessment}
%=============================================================================

\subsection{Configuration 1: SQMS 4-Cavity Array (Incoherent)}

\begin{tcolorbox}[colback=green!5!white, colframe=green!50!black,
  title=\textbf{RECOMMENDED FIRST EXPERIMENT}]

\textbf{Description}: Four TESLA 9-cell cavities at Fermilab SQMS
facility, each independently measuring $Q$-factor constraints on
$\lamdiff$. Results averaged (incoherent $\sqrt{N}$ improvement).

\textbf{Reach}: $\lamdiff < 1.56\ee{-9}/\sqrt{4} = 7.8\ee{-10}$.
Improvement: $2\times$ over current SQMS bound.

\textbf{Cost}: \$2--4M.
\begin{itemize}[nosep]
\item 4 TESLA cavities (may reuse existing SQMS inventory): \$0--800k
\item SQUID readout upgrade (4 channels): \$200k
\item Cryogenic infrastructure: existing at Fermilab
\item Personnel (3 FTE $\times$ 2 yr): \$1.5M
\item Systematic studies, analysis: \$500k
\end{itemize}

\textbf{Timeline}: 1--2 years from approval. All technology exists.
Fermilab SQMS has the infrastructure, expertise, and cavities.

\textbf{Systematic risks}:
\begin{itemize}[nosep]
\item \textbf{Correlated vibration}: All cavities on the same floor slab.
  Mitigation: vibration isolation, time-staggered measurements.
\item \textbf{Correlated magnetic fields}: External field fluctuations
  affect all cavities simultaneously. Mitigation: mu-metal shielding,
  active field cancellation.
\item \textbf{Statistical}: $\sqrt{4} = 2\times$ is modest. If the current
  SQMS bound has unrecognised systematics, the 4-cavity average will not
  improve beyond them.
\end{itemize}

\textbf{Why this is \#1}: Lowest cost, shortest timeline, uses existing
infrastructure, no new technology required. The $2\times$ improvement is
modest but establishes the methodology for scaling to larger arrays.
Any $\lamdiff$ measurement below $1.56\ee{-9}$ would be the strongest
constraint on DFD EM-psi coupling ever achieved.

\textbf{Decision criterion}: If the 4-cavity result is consistent with
$\lamdiff < 7.8\ee{-10}$ and no systematic floor is detected, proceed
to Phase~II (MEMS hybrid) and begin planning Phase~III (256-cavity).
If a systematic floor is detected above $10^{-9}$, diagnose before scaling.
\end{tcolorbox}

\subsection{Configuration 4: P2+P11 Hybrid with $\mu$g MEMS}

\textbf{Description}: Single re-entrant P11 cavity as psi-field source,
coupled to a $\mu$g-scale ($10^{-6}$~kg) Nb membrane MEMS detector with
SQUID readout.

\textbf{Reach}: $\lamdiff < 4.7\ee{-11}$. Improvement: $33\times$ over
SQMS bound.

\textbf{Cost}: \$3--6M.
\begin{itemize}[nosep]
\item Re-entrant cavity design and fabrication: \$300k
\item MEMS Nb membrane fabrication: \$200k (commercial foundry)
\item SQUID readout for MEMS displacement: \$100k
\item Cryogenic integration (MEMS at 2~K in SRF environment): \$500k
\item Personnel (4 FTE $\times$ 3 yr): \$3M
\item Systematic characterisation: \$500k
\end{itemize}

\textbf{Timeline}: 2--3 years. The $\mu$g Nb membrane is near-term
technology (demonstrated in LIGO prototype work). The main development
is cryogenic MEMS--SRF integration.

\textbf{Systematic risks}:
\begin{itemize}[nosep]
\item \textbf{Mechanical $Q$ at 2~K}: The MEMS $Q_{\rm mech}$ may degrade
  in the SRF magnetic environment. Required: $Q_{\rm mech} > 10^4$;
  assumed: $10^6$.
\item \textbf{Electromagnetic crosstalk}: The SRF cavity's stray fields
  may drive the MEMS resonator directly (EM background, not psi signal).
  Requires careful shielding and null tests.
\item \textbf{Vibration coupling}: Seismic vibrations couple to the MEMS
  at all frequencies. Requires vibration isolation to $<10^{-15}$~m/$\sqrt{\rm Hz}$.
\end{itemize}

\textbf{Why this is Phase~II}: It represents the first experiment that
uses a fundamentally different detection principle (MEMS displacement)
from the Q-factor constraint method. If it achieves $\lamdiff < 10^{-10}$,
it would be the first measurement sensitive to DFD effects $>30\times$
below the SQMS bound.

\subsection{Configuration 11: LIGO 6th Channel}

\textbf{Description}: Search for a psi-field signal in LIGO/Virgo
strain data that is not accounted for by GR gravitational waves.
The ``6th channel'' refers to the psi-mode (scalar breathing mode)
that DFD predicts in addition to the 5 standard GW polarisations.

\textbf{Reach}: $\lamdiff < 6\ee{-19}$ --- by far the deepest reach
of any configuration.

\textbf{Cost}: \$0.5--2M (data analysis only; no hardware).

\textbf{Timeline}: 2--3 years (data exists; analysis methodology needed).

\textbf{Systematic risks}:
\begin{itemize}[nosep]
\item \textbf{EXTREME}: The psi-signal must be extracted from LIGO strain
  data dominated by GR GW signals, instrumental noise, and glitches.
  No methodology for this extraction exists.
\item \textbf{Model dependence}: The predicted psi-waveform depends on the
  DFD theory being correct. A null result constrains $\lamdiff$ only if
  the waveform prediction is trusted.
\item \textbf{Degeneracy}: A scalar breathing mode could be produced by
  other modified gravity theories (Brans-Dicke, $f(R)$, etc.), making
  a detection ambiguous.
\end{itemize}

\textbf{Why this is ranked \#11 despite best reach}: The systematic risks
are so severe that the nominal $6\ee{-19}$ reach is likely
unachievable in practice. The LIGO collaboration has searched for scalar
modes and found none, but their searches were not optimised for the
DFD psi-waveform. A dedicated search might improve constraints by
orders of magnitude, but the path from ``LIGO data exists'' to
``$\lamdiff < 6\ee{-19}$'' is extremely uncertain.

\textbf{Recommendation}: Pursue as a low-cost parallel track alongside
the primary roadmap. If a DFD-specific LIGO analysis pipeline can be
developed, it could provide the strongest constraint at minimal hardware
cost.

\subsection{Configurations 9--10: ESS Lund}

\textbf{Description}: The European Spallation Source (ESS) in Lund,
Sweden, produces $\sim 5$~MW of spallation neutrons. The DFD psi-field
from the neutron target mass modulation can be measured via
cross-correlation with a remote SRF cavity detector.

\textbf{Reach}: $10^{-12}$--$10^{-14}$ (realistic to optimistic).

\textbf{Cost}: \$5--15M (detector construction + beam time).

\textbf{Timeline}: 3--5 years (requires ESS beam time allocation,
which is competitive).

\textbf{Systematic risks}:
\begin{itemize}[nosep]
\item Neutron spallation produces enormous EM backgrounds (gamma rays,
  neutron activation). The psi-signal must be extracted via
  cross-correlation at the pulsed source frequency.
\item Vibration from the MW-class target.
\item ESS beam time allocation is competitive and may not prioritise
  a DFD measurement.
\end{itemize}

\textbf{Assessment}: High potential but requires institutional buy-in
from ESS. Best pursued after Phase~I demonstrates improved SQMS bounds.

%=============================================================================
\section{The Three-Phase Detection Roadmap}
\label{sec:roadmap}
%=============================================================================

\begin{tcolorbox}[colback=blue!5!white, colframe=blue!60!black,
  title=\textbf{DEFINITIVE DFD DETECTION ROADMAP}]

\textbf{Phase I: Establish Methodology (2027--2028)}\\
\textit{Cost: \$2--4M. Timeline: 1--2 yr. Target: $\lamdiff < 7.8\ee{-10}$.}
\begin{itemize}[nosep]
\item Build: SQMS 4-cavity incoherent array at Fermilab.
\item Method: Independent Q-factor monitoring of 4 TESLA cavities,
  statistical combination.
\item Personnel: 3 FTE (2 experimentalists + 1 theorist).
\item Decision gate: If systematic floor detected above $10^{-9}$,
  diagnose before proceeding. If $\lamdiff < 7.8\ee{-10}$ achieved,
  proceed to Phase~II.
\item Parallel track: Begin LIGO 6th-channel data analysis (\$0.5M,
  2--3 postdocs).
\end{itemize}

\textbf{Phase II: New Detection Principle (2028--2030)}\\
\textit{Cost: \$3--6M. Timeline: 2--3 yr. Target: $\lamdiff < 4.7\ee{-11}$.}
\begin{itemize}[nosep]
\item Build: P2+P11 hybrid detector with $\mu$g MEMS proof mass.
\item Method: Re-entrant SRF source cavity + Nb membrane MEMS detector +
  SQUID readout.
\item Key development: Cryogenic MEMS--SRF integration.
\item Personnel: 4 FTE (2 experimentalists + 1 MEMS specialist +
  1 theorist).
\item Decision gate: If $\lamdiff < 10^{-10}$ achieved, this is
  $>15\times$ below SQMS and constitutes a major constraint. Proceed
  to Phase~III.
\item Parallel track: Submit ESS beam-time proposal.
\end{itemize}

\textbf{Phase III: Deep Reach (2030--2035)}\\
\textit{Cost: \$60--120M (array) or \$5--10M (ng MEMS). Timeline: 3--5 yr.}\\
\textit{Target: $\lamdiff < 10^{-12}$.}

Two parallel options:
\begin{itemize}[nosep]
\item \textbf{Option A}: 256-cavity coherent array. Reach: $6.1\ee{-12}$.
  Cost: \$80--120M. Requires phase-coherent psi-beam technology (undemonstrated).
  High infrastructure cost but $\Meff$-independent.
\item \textbf{Option B}: P2+P11 hybrid with ng MEMS. Reach: $1.5\ee{-12}$.
  Cost: \$5--10M. Requires ng-scale MEMS with $Q_{\rm mech} = 10^6$ at 2~K.
  Lower cost, higher technology risk.
\item \textbf{Decision}: Choose based on Phase~II MEMS results. If
  $Q_{\rm mech} > 10^5$ demonstrated at 2~K, Option~B is preferred
  (10$\times$ cheaper). Otherwise, Option~A.
\end{itemize}
\end{tcolorbox}

%=============================================================================
\section{Which Experiment Should Be Built FIRST?}
\label{sec:first_experiment}
%=============================================================================

\begin{tcolorbox}[colback=green!5!white, colframe=green!50!black,
  title=\textbf{ANSWER: SQMS 4-Cavity Incoherent Array at Fermilab}]

\textbf{The first experiment to build is an incoherent 4-cavity SQMS
array at Fermilab SQMS.}

Justification:
\begin{enumerate}
\item \textbf{Lowest cost}: \$2--4M, within single-PI grant range.
\item \textbf{Shortest timeline}: 1--2 years, using existing infrastructure.
\item \textbf{No new technology}: All components (TESLA cavities, SQUID
  readout, cryogenics) are TRL~8--9 at Fermilab.
\item \textbf{Meaningful improvement}: $2\times$ better than current SQMS
  bound. While modest, this is the first step in a systematic programme.
\item \textbf{Systematic diagnostic}: The 4-cavity array probes whether
  the current SQMS bound is statistics-limited (in which case
  $\sqrt{N}$ improvement works) or systematics-limited (in which case
  diagnosis is essential before investing in larger arrays).
\item \textbf{Risk management}: If the 4-cavity result reveals an
  unexpected systematic floor, this is learned at \$3M, not \$100M.
\item \textbf{Institutional feasibility}: Fermilab SQMS programme already
  operates multiple TESLA cavities. This experiment is an incremental
  extension of existing operations, not a new facility.
\end{enumerate}

\textbf{The LIGO 6th-channel analysis should proceed in parallel} as
a low-cost (\$0.5--2M) high-reward secondary track. If successful, it
could leapfrog the entire Phase~I--III programme with $\lamdiff < 10^{-19}$,
but the systematic uncertainties are severe.

\textbf{The P2+P11 MEMS hybrid should NOT be built first} because:
\begin{itemize}[nosep]
\item It requires cryogenic MEMS--SRF integration, which has never been
  demonstrated.
\item If the integration fails, 2--3 years and \$3--6M are spent with
  no physics result.
\item The 4-cavity array provides a guaranteed result (either improved
  bound or systematic diagnostic) within 1--2 years.
\end{itemize}
\end{tcolorbox}

%=============================================================================
\section{Risk Matrix: What Could Go Wrong}
\label{sec:risk_matrix}
%=============================================================================

\begin{table}[H]
\centering
\caption{Risk matrix for the three-phase detection roadmap.}
\begin{tabular}{@{}p{3cm}p{2cm}p{2cm}p{5.5cm}@{}}
\toprule
Risk & Likelihood & Impact & Mitigation \\
\midrule
$\lam = 1$ exactly (no coupling) & Unknown & Programme-ending &
  This is the existential risk. No mitigation possible --- if $\lam = 1$,
  all P10--P11 predictions are void. Phase~I is the cheapest way to
  improve the constraint. \\
\midrule
SQMS bound has unrecognised systematics & Medium & Delays Phase~III &
  Phase~I 4-cavity array is designed to diagnose this. \\
\midrule
MEMS $Q_{\rm mech}$ degrades at 2~K in SRF & Medium & Delays Phase~II &
  Begin MEMS cryogenic characterisation during Phase~I (parallel R\&D). \\
\midrule
256-cavity phase coherence unachievable & Medium--High & Eliminates
  Option~A & Retain Option~B (ng MEMS) as backup. \\
\midrule
LIGO analysis yields no useful constraint & High & Loses parallel
  track & Acceptable; primary roadmap (Phases~I--III) is unaffected. \\
\midrule
ESS beam-time not allocated & Medium & Loses Config.~9--10 &
  ESS is a parallel opportunity, not on the critical path. \\
\midrule
Funding insufficient for Phase~III & Medium & Caps at $\lamdiff \sim
  10^{-11}$ & Phase~II with $\mu$g MEMS still reaches $4.7\ee{-11}$,
  which is a $33\times$ improvement over SQMS. \\
\bottomrule
\end{tabular}
\label{tab:risk_matrix}
\end{table}

%=============================================================================
\section{P11 Summary After W3i6}
\label{sec:p11_summary}
%=============================================================================

\begin{table}[H]
\centering
\caption{P11 W3i6 roadmap summary.}
\begin{tabular}{@{}p{3cm}llllp{2.5cm}@{}}
\toprule
Phase & Config. & $\lamdiff$ & Cost (\$M) & Timeline & Decision Gate \\
\midrule
I & 4-cavity SQMS & $7.8\ee{-10}$ & 2--4 & 2027--28 &
  Systematic floor? \\
(parallel) & LIGO 6th ch. & ($6\ee{-19}$) & 0.5--2 & 2027--29 &
  Pipeline feasible? \\
II & MEMS $\mu$g hybrid & $4.7\ee{-11}$ & 3--6 & 2028--30 &
  $Q_{\rm mech}$ at 2~K? \\
III-A & 256-cavity coh. & $6.1\ee{-12}$ & 80--120 & 2030--35 &
  Phase coherence? \\
III-B & MEMS ng hybrid & $1.5\ee{-12}$ & 5--10 & 2030--34 &
  ng MEMS at $Q = 10^6$? \\
\bottomrule
\end{tabular}
\label{tab:roadmap_summary}
\end{table}

%=============================================================================
%                    PART III: COMBINED FINDINGS
%=============================================================================

\newpage
\part{Combined Findings, Corrections, and Open Questions}

%=============================================================================
\section{Ranked Findings (Combined P8+P11, W3i6)}
\label{sec:ranked_findings}
%=============================================================================

\begin{enumerate}

\item \textbf{[RANK 1 --- ROADMAP] Definitive P11 Detection Roadmap.}
  Three-phase programme: Phase~I SQMS 4-cavity (\$2--4M, 2027--28,
  $7.8\ee{-10}$), Phase~II MEMS hybrid (\$3--6M, 2028--30, $4.7\ee{-11}$),
  Phase~III deep reach (\$5--120M, 2030--35, $10^{-12}$).
  The first experiment to build is the 4-cavity SQMS array at Fermilab.

\item \textbf{[RANK 2 --- KILLER APP] Submarine/underground communication
  is the ONE irreplaceable P8 application.}
  No alternative physics provides bidirectional real-time communication
  through arbitrary matter thickness. However, the transmitter mass gap
  ($10^4\times$) makes it currently infeasible.

\item \textbf{[RANK 3 --- CONFIRMED] Two-component $\psi$ does not change
  any P8 prediction.}
  Secrecy, corridor propagation, link budget, and noise floor are all
  unaffected. P8 communications operates in the $\psigrav$ sector via
  mass modulation.

\item \textbf{[RANK 4 --- STRUCTURAL] The LIGO 6th channel is the
  highest-reward, highest-risk option.}
  At $\lamdiff < 6\ee{-19}$ it dwarfs all other configurations by
  $>10^6\times$ in reach, but the analysis methodology is undemonstrated
  and systematic risks are extreme. Should be pursued as a parallel
  low-cost track.

\item \textbf{[RANK 5 --- CONFIRMED] Passive Superiority Theorem (P2)
  drives the roadmap.}
  All viable detection paths use either $\Meff$-independent observables
  (Q-factor, optical path, cross-correlation) or sub-kg MEMS proof masses.
  No bulk SRF SQUID-active path can beat SQMS. The W3i5 Passive
  Superiority Theorem is the organising principle of the roadmap.

\item \textbf{[RANK 6 --- HONEST] Deep-space psi-comms is definitively
  non-competitive.}
  DSN wins by $370\times$ on Mars rate. Laser comms (DSOC) will extend EM
  advantage further. The P8 deep-space claims should be classified as
  \emph{theoretical demonstrations} of DFD physics, not \emph{technological
  proposals}.

\item \textbf{[RANK 7 --- RISK] Existential risk: $\lam = 1$ exactly.}
  If $\lam = 1$, all P10--P11--P15 predictions collapse and there is no
  EM-psi coupling to detect. The roadmap is designed to constrain $\lamdiff$
  at minimum cost per decade of improvement.

\end{enumerate}

%=============================================================================
\section{Required Corpus Corrections (W3i6)}
\label{sec:corpus_corrections}
%=============================================================================

\begin{enumerate}
\item \textbf{Add to P8 section}: ``The ONE killer application for
  psi-communications is bidirectional submarine/underground communication.
  All other applications are dominated by EM alternatives. Even the
  killer application requires $M_s \gtrsim 10^{10}$~kg (currently
  infeasible).''

\item \textbf{Add to P11 section}: ``Definitive three-phase detection
  roadmap: Phase~I (4-cavity SQMS, \$2--4M, 2027--28), Phase~II
  (MEMS $\mu$g hybrid, \$3--6M, 2028--30), Phase~III (256-cavity or
  ng MEMS, \$5--120M, 2030--35). First experiment: 4-cavity SQMS at
  Fermilab.''

\item \textbf{Update P8 status}: Reclassify all deep-space communications
  claims as ``theoretical demonstrations'' rather than ``technological
  proposals.''

\item \textbf{Confirm}: Two-component $\psi$ decomposition has zero impact
  on P8 link budget, secrecy, corridor propagation, and noise floor.

\item \textbf{Add risk matrix}: The existential risk $\lam = 1$ should be
  flagged prominently in the Master Corpus as the overarching uncertainty
  for all P10--P11--P15 predictions.
\end{enumerate}

%=============================================================================
\section{Open Questions for Iteration 7}
\label{sec:open_questions}
%=============================================================================

\begin{enumerate}
\item \textbf{P8 transmitter mass reduction}: Can P2 passive amplification
  or coherent mass-modulation techniques reduce the required $M_s$ for
  submarine comms? If the Passive Superiority Theorem applies to
  \emph{transmission} (not just detection), there may be a path to
  smaller source masses through resonant enhancement.

\item \textbf{P11 MEMS cryogenic characterisation}: What is the measured
  $Q_{\rm mech}$ of Nb membranes and MEMS beams at 2~K in a 1.3~GHz
  SRF magnetic field environment? This is the single most important
  technological input for the Phase~II decision gate.

\item \textbf{P11 LIGO analysis pipeline}: Can a DFD-specific waveform
  template be constructed for the scalar breathing mode? What is the
  sensitivity of existing LIGO O3/O4 data to this template?

\item \textbf{P11 phase coherence}: What is the theoretical and practical
  limit to maintaining psi-wave phase coherence across $N$ spatially
  separated SRF cavities? This determines whether the $1/N$ (coherent)
  or $1/\sqrt{N}$ (incoherent) scaling applies to the 256-cavity array.

\item \textbf{Two-component $\psi$ and P15}: Under the two-component
  decomposition, P15 (psi-field generation) operates by converting EM
  energy into $\dpsiEM$. Does the two-component picture change the P15
  generation efficiency? (P15 was killed in W3i5 by the $\Meff$ correction;
  does the two-component picture offer any recovery?)
\end{enumerate}

\end{document}
